% =====================================================================
%  Relatório — INF05010 Otimização Combinatória (UFRGS)
%  Problema: Os Melhores Amigos (OMA) — Meta-heurística: Busca Tabu
%
%  Formatação segundo as normas oficiais SOBRAPO/SBPO (Artigos Completos):
%    A4, fonte Times New Roman 11, coluna única, espaçamento simples,
%    margens 3,3 / 2,5 / 2,9 cm, no máximo 12 páginas (ver docs/sobrapo_norms.md).
%  O ideal é substituir este preâmbulo pelo template LaTeX oficial PT-BR da
%  SOBRAPO; este reproduz as mesmas normas e compila sem .sty externo.
% =====================================================================
\documentclass[a4paper,11pt]{article}

\usepackage[T1]{fontenc}
\usepackage[utf8]{inputenc}
\usepackage[brazil]{babel}
\usepackage{mathptmx}            % Times New Roman
\usepackage{amsmath,amssymb}
\usepackage[top=3.3cm,bottom=2.5cm,left=2.9cm,right=2.9cm]{geometry}
\usepackage{graphicx}
\usepackage{booktabs}
\usepackage{array}
\usepackage{siunitx}             % colunas numéricas alinhadas (S)
\usepackage{algorithm}
\usepackage{algpseudocode}
\usepackage{url}
\usepackage{cite}

\setlength{\parindent}{1.25em}

\title{\textbf{Os Melhores Amigos}}

% Submissão ao SBPO seria ANÔNIMA (sem autores); como este é um trabalho de
% disciplina, os autores são incluídos na versão de entrega.
\author{Gabriel Stoffel \quad Rodrigo Feldens \\
  \small INF05010 --- Otimização Combinatória \\
  \small Instituto de Informática, Universidade Federal do Rio Grande do Sul (UFRGS)}
\date{\today}

\begin{document}
\maketitle

% ---------------------------------------------------------------------
\begin{abstract}
\noindent\textbf{Resumo.}
Este trabalho aborda o problema \emph{Os Melhores Amigos} (OMA): dado um conjunto de
$n$ pessoas e um valor de afinidade não negativo para cada par, deseja-se selecionar
exatamente $m$ pessoas de modo a maximizar a soma das afinidades entre os pares
escolhidos. O problema é tratado por duas abordagens complementares: (i) uma
formulação como programa linear inteiro, resolvida por um solver genérico
(HiGHS/CPLEX) que fornece a referência ótima por instância; e (ii) uma meta-heurística
de \emph{Busca Tabu}, com construção inicial gulosa, vizinhança por troca e avaliação
incremental do objetivo. O relatório descreve o modelo matemático, os componentes da
Busca Tabu e o protocolo experimental, e (nas próximas versões) reportará a análise de
sensibilidade dos parâmetros e a comparação entre a meta-heurística e o método exato
sobre as 10 instâncias de teste.

\medskip
\noindent\textbf{Palavras-chave:} Busca Tabu; Otimização Combinatória; Programação Linear Inteira.
\end{abstract}

% ---------------------------------------------------------------------
\section{Introdução}
\label{sec:intro}

O problema \emph{Os Melhores Amigos} (OMA) consiste em, dado um conjunto de
$n$ pessoas $P = [n]$ e um valor de afinidade não negativo $a_{ij} \geq 0$ para cada
par $\{i,j\}$ com $i < j$, escolher um subconjunto $S \subseteq P$ com \textbf{exatamente}
$m$ pessoas que maximize a afinidade total do grupo,
%
\begin{equation}
  A(S) = \sum_{\{i,j\} \subseteq S} a_{ij}.
  \label{eq:obj}
\end{equation}
%
Trata-se de um problema de seleção combinatória: o número de subconjuntos candidatos é
$\binom{n}{m}$, o que torna a enumeração exaustiva impraticável já para instâncias de
tamanho moderado. Na literatura de Pesquisa Operacional, o OMA corresponde ao
\textbf{Problema da Máxima Diversidade} (\emph{Maximum Diversity Problem}, MDP) em sua
variante \emph{MaxSum}, no qual se seleciona um subconjunto de cardinalidade fixa que
maximiza a soma das afinidades (ou dist\^ancias) entre os elementos
escolhidos~\cite{kuo1993analyzing,marti2013heuristics,marti2022review}; o problema é
NP-difícil~\cite{kuo1993analyzing}.

Neste trabalho o problema é atacado de duas formas complementares. A primeira é uma
\textbf{formulação como programa linear inteiro} (Seção~\ref{sec:formulacao}), resolvida
por um solver genérico de software livre, que fornece o valor ótimo --- ou o melhor
limite encontrado dentro de um tempo-limite --- e serve de \textbf{referência} contra a
qual a heurística é avaliada. A segunda é uma \textbf{meta-heurística de Busca Tabu}
(Seção~\ref{sec:solucao}), capaz de obter boas soluções em tempo reduzido. O objetivo
central do estudo é medir quão perto a Busca Tabu chega da referência ótima e a que
custo de tempo, além de analisar a sensibilidade do método aos seus parâmetros.

A avaliação utiliza um conjunto fixo de \textbf{10 instâncias} ($\texttt{oma01}$ a
$\texttt{oma10}$): as cinco primeiras com $n = 125$ e $m = 12$, e as cinco últimas com
$n = 150$ e $m = 15$. Os melhores valores conhecidos (BKV) dessas instâncias estão
reproduzidos na Tabela~\ref{tab:bkv} e são usados como referência quando o solver não
prova a otimalidade.

\begin{table}[t]
  \centering
  \caption{Melhores valores conhecidos (BKV) das instâncias de teste.}
  \label{tab:bkv}
  \begin{tabular}{l S[table-format=3.0] l S[table-format=3.0]}
    \toprule
    {Instância} & {BKV} & {Instância} & {BKV} \\
    \midrule
    oma01 & 472 & oma06 & 719 \\
    oma02 & 474 & oma07 & 724 \\
    oma03 & 470 & oma08 & 732 \\
    oma04 & 470 & oma09 & 733 \\
    oma05 & 474 & oma10 & 721 \\
    \bottomrule
  \end{tabular}
\end{table}

O restante do relatório está organizado da seguinte forma. A Seção~\ref{sec:formulacao}
apresenta a formulação inteira do problema. A Seção~\ref{sec:solucao} descreve a
meta-heurística de Busca Tabu e seus parâmetros. A Seção~\ref{sec:resultados} reporta os
resultados experimentais e a análise, e a Seção~\ref{sec:conclusao} traz as conclusões.

% ---------------------------------------------------------------------
\section{Formulação como Programa Inteiro}
\label{sec:formulacao}

O OMA pode ser modelado como um programa linear inteiro binário. Para cada pessoa
$i \in [n]$ define-se a variável de decisão $x_i \in \{0,1\}$, com $x_i = 1$ se e somente
se a pessoa $i$ é selecionada. Como o objetivo~\eqref{eq:obj} envolve o produto
$x_i\,x_j$ (a afinidade de um par só conta quando ambas as pessoas são escolhidas),
introduz-se, para cada par $i < j$, uma variável auxiliar $z_{ij} \in \{0,1\}$ que
\emph{lineariza} esse produto: deseja-se $z_{ij} = 1$ apenas quando $x_i = x_j = 1$. Essa
é a linearização clássica de produtos de variáveis binárias~\cite{glover1974converting},
aplicada ao MDP por Kuo, Glover e Dhir~\cite{kuo1993analyzing}.

O modelo completo é:
%
\begin{align}
  \max \quad & \sum_{\substack{i,j \in [n] \\ i < j}} a_{ij}\, z_{ij} \label{eq:ilp-obj}\\
  \text{s.a.} \quad
    & \sum_{i=1}^{n} x_i = m \label{eq:card}\\
    & z_{ij} \leq x_i & & \forall\, i,j \in [n],\; i < j \label{eq:lin1}\\
    & z_{ij} \leq x_j & & \forall\, i,j \in [n],\; i < j \label{eq:lin2}\\
    & z_{ij} \geq x_i + x_j - 1 & & \forall\, i,j \in [n],\; i < j \label{eq:lin3}\\
    & x_i \in \{0,1\} & & \forall\, i \in [n] \label{eq:intx}\\
    & z_{ij} \in \{0,1\} & & \forall\, i,j \in [n],\; i < j \label{eq:intz}
\end{align}

\paragraph{Explicação das restrições.}
A restrição de \textbf{cardinalidade}~\eqref{eq:card} garante que exatamente $m$ pessoas
sejam selecionadas. As restrições~\eqref{eq:lin1} e~\eqref{eq:lin2} forçam $z_{ij} = 0$
sempre que pelo menos uma das pessoas do par não estiver na solução, isto é, $z_{ij}$ só
pode valer $1$ se $x_i = 1$ e $x_j = 1$. A restrição~\eqref{eq:lin3} atua no sentido
oposto: quando $x_i = x_j = 1$, ela impõe $z_{ij} \geq 1$ e, portanto, $z_{ij} = 1$,
assegurando que a afinidade do par seja efetivamente contabilizada no
objetivo~\eqref{eq:ilp-obj}. As restrições~\eqref{eq:intx} e~\eqref{eq:intz} impõem a
integralidade das variáveis. Como o objetivo é de maximização e todos os $a_{ij}$ são não
negativos, as desigualdades~\eqref{eq:lin1}--\eqref{eq:lin2} bastam para limitar
$z_{ij}$ por cima; a desigualdade~\eqref{eq:lin3} é necessária para impedir que o solver
deixe $z_{ij} = 0$ com ambas as pessoas selecionadas.

\paragraph{Tamanho do modelo.}
O modelo usa $n$ variáveis $x_i$ e $\binom{n}{2}$ variáveis $z_{ij}$, além de uma
restrição de cardinalidade e três restrições de linearização por par, totalizando
$3\binom{n}{2} + 1$ restrições. O crescimento quadrático no número de pares ($z_{ij}$ e
restrições) é o principal fator que limita a resolução exata em instâncias maiores e
motiva o uso da meta-heurística.

% ---------------------------------------------------------------------
\section{Descrição da Solução --- Busca Tabu}
\label{sec:solucao}

A meta-heurística escolhida é a \textbf{Busca Tabu} (BT)~\cite{glover1989tabu,glover1997tabu},
uma busca local que escapa de ótimos locais mantendo uma memória de curto prazo (a
\emph{lista tabu}) que proíbe temporariamente o desfazimento de movimentos recentes
\cite{hertz1995tutorial,gendreau2003intro,ritt2024inf05010}. A aplicação de Busca Tabu e
de procedimentos gulosos/GRASP ao MDP é bem estabelecida na
literatura~\cite{duarte2007tabu}. A seguir descrevem-se seus componentes.

\subsection{Representação da solução}
Uma solução é um subconjunto $S \subseteq P$ com cardinalidade $|S| = m$. Todas as
soluções manipuladas pela busca são \emph{factíveis} por construção: a cardinalidade
$m$ é preservada em cada passo.

\subsection{Solução inicial}
A solução inicial é construída por uma heurística \textbf{gulosa}~\cite{feo1995greedy} que,
partindo do conjunto vazio, adiciona repetidamente a pessoa que mais aumenta a afinidade
total até atingir $m$ pessoas:
%
\begin{equation*}
  \text{ganho}(p) = \sum_{i \in S} a_{ip}, \qquad
  p^\star = \arg\max_{p \notin S} \text{ganho}(p).
\end{equation*}
%
Para evitar recalcular a soma a cada iteração, mantém-se um vetor de ganhos incrementais
atualizado a cada inserção. Como alternativa de comparação (bloco experimental A5),
considera-se também uma \textbf{construção aleatória}, que amostra $m$ pessoas
uniformemente, permitindo avaliar o ganho efetivo da partida gulosa.

\subsection{Vizinhança}
A partir de uma solução $S$, um vizinho $S'$ é obtido por uma \textbf{troca} (\emph{swap}):
remove-se uma pessoa $p \in S$ e insere-se uma pessoa $p' \notin S$,
%
\begin{equation*}
  S' = (S \setminus \{p\}) \cup \{p'\}, \qquad p \in S,\; p' \notin S.
\end{equation*}
%
A vizinhança tem tamanho $m\,(n-m)$ e preserva a cardinalidade $|S'| = m$.

\subsection{Função objetivo e avaliação incremental (delta)}
Recalcular $A(S)$ inteira a cada vizinho é custoso. Em vez disso, avalia-se o
\textbf{delta} de uma troca de $p$ por $p'$ em $O(m)$:
%
\begin{equation}
  \Delta = A(S') - A(S)
         = \sum_{i \in S \setminus \{p\}} a_{p'i}
         - \sum_{i \in S \setminus \{p\}} a_{pi}.
  \label{eq:delta}
\end{equation}
%
A cada iteração escolhe-se a troca não tabu de maior $\Delta$ (estratégia de melhor
vizinho).

\subsection{Lista tabu e \emph{tenure}}
A lista tabu $T$ armazena as pessoas recentemente removidas da solução, impedindo sua
reinserção imediata por um número fixo de iterações chamado \emph{tenure}. Esse mecanismo
evita ciclos curtos e força a busca a explorar regiões diferentes do espaço de soluções.

\subsection{Critério de aspiração}
Um movimento marcado como tabu é, ainda assim, permitido se ele produz uma solução melhor
do que a \textbf{melhor solução global} encontrada até o momento. Esse critério de
aspiração evita descartar movimentos comprovadamente vantajosos apenas por estarem na
lista tabu.

\subsection{Diversificação}
Para escapar de ótimos locais, aplica-se periodicamente uma \textbf{perturbação aleatória}
da solução corrente (a cada \texttt{diversify\_freq} iterações), reintroduzindo
diversidade na busca. Quando \texttt{diversify\_freq} $= 0$, a diversificação é desligada.

\subsection{Critérios de parada}
A busca termina quando (i) atinge um número máximo de iterações (\texttt{max\_iter}) ou
(ii) acumula um número máximo de iterações consecutivas sem melhora na melhor solução
(\texttt{max\_no\_improve}).

\subsection{Pseudocódigo}
O Algoritmo~\ref{alg:bt} resume o procedimento.

\begin{algorithm}[t]
  \caption{Busca Tabu para o OMA}
  \label{alg:bt}
  \begin{algorithmic}[1]
    \State $S \gets \textsc{ConstrucaoInicial}()$ \Comment{gulosa ou aleatória}
    \State $S^\star \gets S$ \Comment{melhor solução global}
    \State $T \gets \emptyset$; \; $iter \gets 0$; \; $semMelhora \gets 0$
    \While{$iter < \texttt{max\_iter}$ \textbf{and} $semMelhora < \texttt{max\_no\_improve}$}
      \State escolha a troca $(p, p')$ que maximiza $\Delta$~\eqref{eq:delta}, não tabu
        \textbf{ou} que satisfaça a aspiração
      \State $S \gets (S \setminus \{p\}) \cup \{p'\}$
      \State marque $p$ como tabu por \emph{tenure} iterações em $T$
      \If{$A(S) > A(S^\star)$}
        \State $S^\star \gets S$; \; $semMelhora \gets 0$
      \Else
        \State $semMelhora \gets semMelhora + 1$
      \EndIf
      \If{\texttt{diversify\_freq} $> 0$ \textbf{and} $iter \bmod \texttt{diversify\_freq} = 0$}
        \State $S \gets \textsc{Perturbacao}(S)$
      \EndIf
      \State $iter \gets iter + 1$
    \EndWhile
    \State \Return $S^\star$
  \end{algorithmic}
\end{algorithm}

\subsection{Parâmetros}
A Tabela~\ref{tab:params} lista os parâmetros do método e seus papéis. Os valores de
referência (\emph{baseline}) usados como ponto central das varreduras de sensibilidade
são \texttt{tenure}~$=10$, \texttt{max\_no\_improve}~$=100$, \texttt{max\_iter}~$=1000$,
\texttt{diversify\_freq}~$=50$ e construção inicial gulosa.

\begin{table}[t]
  \centering
  \caption{Parâmetros da Busca Tabu.}
  \label{tab:params}
  \begin{tabular}{l l}
    \toprule
    Parâmetro & Papel \\
    \midrule
    \texttt{tenure}          & duração (em iterações) da proibição de reinserção \\
    \texttt{max\_no\_improve} & iterações sem melhora toleradas antes de parar \\
    \texttt{max\_iter}        & número máximo de iterações \\
    \texttt{diversify\_freq}  & frequência da perturbação ($0$ = desligada) \\
    \texttt{init}             & construção inicial: gulosa ou aleatória \\
    \bottomrule
  \end{tabular}
\end{table}

% ---------------------------------------------------------------------
\section{Resultados Obtidos e Análise}
\label{sec:resultados}

% TODO: preencher após rodar os experimentos (varredura M1/M2/M3 + solver) e
% gerar figuras/tabelas com `uv run oma-report --results-dir results --out-dir report`.
% Esta seção espelhará a estrutura do relatório de referência (docs/reference_report.md):
%   - Análise de sensibilidade dos parâmetros (blocos A1--A4), figuras report/figures/A*.png
%   - Comparação da construção inicial gulosa vs. aleatória (bloco A5)
%   - Tabelas por instância (baseline/tunada/degenerada vs. solver vs. BKV), report/tables/T*.{csv,md}
%   - Comparação exato vs. exato (HiGHS vs. CPLEX)
\textit{Seção a ser preenchida após a execução dos experimentos.}

% ---------------------------------------------------------------------
\section{Conclusão}
\label{sec:conclusao}

% TODO: preencher após a análise dos resultados.
\textit{Seção a ser preenchida após a análise dos resultados.}

% ---------------------------------------------------------------------
\bibliographystyle{plain}
\bibliography{referencias}

\end{document}
